\setlength{\absparsep}{18pt}

\begin{resumo}

Este trabalho teve como objetivo desenvolver um sistema web de apoio à avaliação psicopedagógica na Educação Infantil que reunisse, em um mesmo ambiente, o cadastro de aprendentes, o gerenciamento de instrumentos avaliativos, o registro das aplicações e a elaboração de relatórios. Trata-se de uma pesquisa aplicada, conduzida segundo a abordagem da \textit{Design Science Research}, na qual os requisitos foram definidos a partir de revisão bibliográfica, de análise documental de instrumentos psicopedagógicos e de entrevista semiestruturada com uma psicopedagoga, e o sistema foi construído de forma incremental. A aplicação adota arquitetura em três camadas, com interface em React, API REST em Node.js e Express e banco de dados MySQL acessado por meio do Prisma. Entre os recursos implementados, destacam-se o controle de acesso por perfil e por titularidade dos registros, a separação entre a administração do sistema e o acesso aos dados clínicos, em respeito ao sigilo profissional, o versionamento automático dos instrumentos, o registro de auditoria e a emissão de relatórios em PDF. Os 147 testes automatizados foram aprovados, com cobertura de 83,06\% das linhas do \textit{backend}, e a verificação automatizada de acessibilidade não apontou violações dos critérios A e AA das WCAG 2.2 após as correções realizadas. Cinco profissionais da psicopedagogia avaliaram o sistema por meio de questionário e atribuíram média geral de 4,43 em uma escala de 1 a 5, sendo a geração de relatórios o aspecto com menor média (3,8). Dos quinze requisitos funcionais, treze foram plenamente atendidos e dois, parcialmente. Conclui-se que o sistema atende ao objetivo proposto e foi percebido como adequado à prática profissional. Permanece, contudo, em estágio de protótipo, uma vez que a ausência de criptografia em repouso, de política de retenção de dados e de interoperabilidade, somada à amostra reduzida da avaliação, impede seu uso com dados reais sem ajustes adicionais.

\noindent
\textbf{Palavras-chave}: psicopedagogia; avaliação psicopedagógica; educação infantil; sistema web; tecnologias educacionais.

\end{resumo}